% ---
% title: A random portfolio
% tags: finance, equally-likely, complement
% ---

\begin{problem}
	\begin{context}
		A stock \emph{index} is a list of stocks whose prices are tracked
		together. In an \emph{equally weighted} portfolio the same amount of
		money goes into each chosen stock, so only the \emph{set} of stocks
		matters. Investors watch \emph{sector concentration}: a portfolio with
		many stocks from one sector (for example technology) is exposed to the
		risks of that sector.
	\end{context}
	An index has 30 stocks, 8 of them technology stocks. An investor picks 5
	of the 30 stocks completely at random for an equally weighted portfolio.
	Find the probability that the portfolio contains (a) at least one
	technology stock; (b) exactly two technology stocks; (c) no technology
	stock.
\end{problem}

\begin{solution}
	\idea ``Completely at random'' means that every set of 5 stocks is equally
	likely. So every probability is (number of favorable sets) divided by
	(number of all sets). For ``at least one'' use the complement: ``none''.

	\step{Sample space} An outcome is a set of 5 stocks out of 30. The order of
	picking does not matter (the portfolio is equally weighted), so
	\[ |\Om|=\binom{30}{5}=142\,506. \]
	All these outcomes are equally likely.

	\step{(c)} We start with (c), because (a) builds on it. ``No technology
	stock'' means that all 5 stocks come from the $30-8=22$ other stocks. The
	number of such sets is $\binom{22}{5}=26\,334$, so
	\[ P(\text{no tech})=\frac{\binom{22}{5}}{\binom{30}{5}}
		=\frac{26\,334}{142\,506}=\frac{209}{1131}\approx0.1848. \]

	\step{(a)} Let $E$ = ``at least one technology stock''. Its complement
	$E^c$ is ``no technology stock'', which we found in (c). Counting $E$
	directly would need the cases 1, 2, 3, 4, 5 tech stocks; the complement is
	one single case. Since $P(E)=1-P(E^c)$,
	\[ P(E)=1-\frac{26\,334}{142\,506}=\frac{116\,172}{142\,506}=\frac{922}{1131}
		\approx0.8152. \]

	\step{(b)} ``Exactly two technology stocks'' means: 2 stocks from the 8
	tech stocks \emph{and} 3 stocks from the 22 others. Any choice of the 2
	tech stocks can be combined with any choice of the 3 other stocks, so by
	the multiplication principle the
	number of such sets is
	\[ \binom{8}{2}\binom{22}{3}=28\cdot1540=43\,120. \]
	Hence
	\[ P(\text{exactly 2 tech})=\frac{\binom{8}{2}\binom{22}{3}}{\binom{30}{5}}
		=\frac{43\,120}{142\,506}=\frac{3080}{10\,179}\approx0.3026. \]

	\begin{answer}
		(a) $\approx0.8152$ \quad (b) $\approx0.3026$ \quad (c) $\approx0.1848$
	\end{answer}
	\begin{remark}
		The counts are ordinary combinations; the only new step is dividing by
		$|\Om|$, which is allowed because all sets of 5 stocks are equally
		likely.

		\textbf{Common mistake:} in (b) writing only $\binom{8}{2}$ in the
		numerator. The other 3 stocks must be \emph{non}-tech stocks; otherwise
		the portfolio could contain more than two tech stocks.
	\end{remark}
\end{solution}
